% Bibliografía para 2_Antecedentes (solo claves citadas: 27)
\begin{thebibliography}{99}
\bibitem{Abdelhafiz2019DeepCNN}
    Abdelhafiz, D., Yang, C., Ammar, R., et al. (2019). Deep convolutional neural networks for mammography:
    advances, challenges and applications. \textit{BMC Bioinformatics}, 20(Suppl 11), 281.
    \url{https://doi.org/10.1186/s12859-019-2823-4}

  % --- Referencia 25: Deep Learning to Improve Breast Cancer Detection on Screening Mammography ---

\bibitem{Anez2025AIPipeline}
    Añez, D., Conti, G., Uriarte, J. J., Serrano-Olmedo, J.-J., Martínez-Murillo, R., \& Casanova-Carvajal, O. (2025).
    Artificial intelligence pipeline for mammography-based breast cancer detection: An integrated systematic review and
    large-scale experimental validation. \textit{Medicina}, 61(12), 2237.
    \url{https://doi.org/10.3390/medicina61122237}

  % --- Referencia 24: Deep Convolutional Neural Networks for Mammography ---

\bibitem{Avci2023MammographyEnhancement}
    Avcı, H., \& Karakaya, J. (2023). A novel medical image enhancement algorithm for breast cancer
    detection on mammography images using machine learning. \textit{Diagnostics}, 13(3), 348.
    \url{https://doi.org/10.3390/diagnostics13030348}

  % --- Referencia 23: AI Pipeline for Mammography Breast Cancer Detection - PRISMA Review ---

\bibitem{Azevedo2022QuantumTransfer}
    Azevedo, V., Silva, C., \& Dutra, I. (2022). Quantum transfer learning for breast
    cancer detection. \textit{Quantum Machine Intelligence}, 4(1), 5.
    \url{https://doi.org/10.1007/s42484-022-00062-4}

  % --- Referencia 17: Quantum Kernel Methods for Marketing Analytics ---

\bibitem{Baccouche2022ResidualNN}
    Baccouche, A., Garcia-Zapirain, B., \& Elmaghraby, A. S. (2022). An integrated framework for breast mass
    classification and diagnosis using stacked ensemble of residual neural networks. \textit{Scientific Reports}, 12, 12259.
    \url{https://doi.org/10.1038/s41598-022-15632-6}

  % --- Referencia 27: Deep Learning Classification of Breast Microcalcifications ---

\bibitem{Bowles2024Benchmarking}
    Bowles, J., Ahmed, S., \& Schuld, M. (2024).
    Better than classical? The subtle art of benchmarking quantum machine learning models.
    \textit{arXiv preprint} arXiv:2403.07059.
    \url{https://doi.org/10.48550/arXiv.2403.07059}

\bibitem{Canatar2023Bandwidth}
    Canatar, A., Peters, E., Pehlevan, C., Wild, S. M., \& Shaydulin, R. (2022).
    Bandwidth enables generalization in quantum kernel models.
    \textit{arXiv preprint} arXiv:2206.06686.
    \url{https://doi.org/10.48550/arXiv.2206.06686}

\bibitem{ElHafidi2025}
    El Hafidi, M. Y., Toufah, A., \& Kadim, M. A. (2025).
    Investigating quantum feature maps in quantum support vector machines for lung cancer classification.
    \textit{arXiv preprint} arXiv:2506.03272.
    \url{https://doi.org/10.48550/arXiv.2506.03272}

\bibitem{Glick2024Covariant}
    Glick, J. R., Gujarati, T. P., C{\'{o}}rcoles, A. D., Kim, Y., Kandala, A., Gambetta, J. M., \& Temme, K. (2024).
    Covariant quantum kernels for data with group structure.
    \textit{Nature Physics}, 20(3), 479--483.
    \url{https://doi.org/10.1038/s41567-023-02340-9}

\bibitem{Huang2022Experiments}
    Huang, H.-Y., Broughton, M., Cotler, J., Chen, S., Li, J., Mohseni, M., Neven, H., Babbush, R., Kueng, R., Preskill, J., \& McClean, J. R. (2022).
    Quantum advantage in learning from experiments.
    \textit{Science}, 376(6598), 1182--1186.
    \url{https://doi.org/10.1126/science.abn7293}

\bibitem{Kubler2021Inductive}
    K\"ubler, J. M., Buchholz, S., \& Sch\"olkopf, B. (2021).
    The inductive bias of quantum kernels.
    En \textit{Advances in Neural Information Processing Systems 34 (NeurIPS 2021)}.
    \url{https://doi.org/10.48550/arXiv.2106.03747}

\bibitem{Le2011FRQI}
    Le, P. Q., Dong, F., \& Hirota, K. (2011).
    A flexible representation of quantum images for polynomial preparation, image compression, and processing operations.
    \textit{Quantum Information Processing}, 10(1), 63--84.
    \url{https://doi.org/10.1007/s11128-010-0177-y}

\bibitem{Liu2021Speedup}
    Liu, Y., Arunachalam, S., \& Temme, K. (2021).
    A rigorous and robust quantum speed-up in supervised machine learning.
    \textit{Nature Physics}, 17(9), 1013--1017.
    \url{https://doi.org/10.1038/s41567-021-01287-z}

\bibitem{Long2025HybridQCQ}
    Long, C., Huang, M., Ye, X., et al. (2025). Hybrid quantum-classical-quantum convolutional neural networks.
    \textit{Scientific Reports}, 15, 31780.
    \url{https://doi.org/10.1038/s41598-025-13417-1}

  % --- Referencia 30: Quantum Machine and Deep Learning for Medical Image Classification ---

\bibitem{Mashekova2025AIBreastCancer}
    Mashekova, A., Zhao, M. Y., Zarikas, V., Mukhmetov, O., Aidossov, N., Ng, E. Y. K., \& Wang, D. (2025).
    Review of artificial intelligence techniques for breast cancer detection with different modalities:
    Mammography, ultrasound, and thermography images. \textit{Bioengineering}, 12(10), 1110.
    \url{https://doi.org/10.3390/bioengineering12101110}

  % --- Referencia 22: Medical Image Enhancement for Breast Cancer Detection ---

\bibitem{Parigi2026QImR}
    Parigi, M., Khosrojerdi, M., Caruso, F., \& Banchi, L. (2026).
    Analysis of quantum image representations for supervised classification.
    \textit{AVS Quantum Science}, 8(1), 013801.
    \url{https://doi.org/10.1116/5.0296376}

\bibitem{Schuld2019QuantumML}
    Schuld, M., \& Killoran, N. (2019). Quantum machine learning in feature {H}ilbert spaces.
    \textit{Physical Review Letters}, 122(4), 040504.
    \url{https://doi.org/10.1103/PhysRevLett.122.040504}

  % --- Referencia 29: Hybrid Quantum-Classical-Quantum Convolutional Neural Networks ---

\bibitem{SchuldKilloran2022Advantage}
    Schuld, M., \& Killoran, N. (2022).
    Is quantum advantage the right goal for quantum machine learning?
    \textit{PRX Quantum}, 3(3), 030101.
    \url{https://doi.org/10.1103/PRXQuantum.3.030101}

\bibitem{Shaydulin2022Bandwidth}
    Shaydulin, R., \& Wild, S. M. (2022).
    Importance of kernel bandwidth in quantum machine learning.
    \textit{Physical Review A}, 106(4), 042407.
    \url{https://doi.org/10.1103/PhysRevA.106.042407}

\bibitem{Shen2019DeepLearning}
    Shen, L., Margolies, L. R., Rothstein, J. H., et al. (2019). Deep learning to improve breast cancer
    detection on screening mammography. \textit{Scientific Reports}, 9, 12495.
    \url{https://doi.org/10.1038/s41598-019-48995-4}

  % --- Referencia 26: Integrated Framework for Breast Mass Classification Using Residual Neural Networks ---

\bibitem{Shia2025MicrocalcClassification}
    Shia, W. C. (2025). Deep learning-based classification of benign and malignant breast microcalcifications in
    mammography. \textit{Scientific Reports}, 15, 42768.
    \url{https://doi.org/10.1038/s41598-025-27149-9}

  % --- Referencia 28: Quantum Machine Learning in Feature Hilbert Spaces ---

\bibitem{TandFOnline2021}
    Author(s). (2021). [Article title]. \textit{AI Magazine}, 42(4), [pages].
    \url{https://doi.org/10.1080/08839514.2021.2001177}

  % --- Referencia 20: Machine Learning and Deep Learning in Mammography ---

\bibitem{Thanasilp2024}
    Thanasilp, S., Wang, S., Cerezo, M., \& Holmes, Z. (2024).
    Exponential concentration in quantum kernel methods.
    \textit{Nature Communications}, 15(1), 5200.
    \url{https://doi.org/10.1038/s41467-024-49287-w}

\bibitem{Xiang2024QuantumCNN}
    Xiang, Q., Li, D., Hu, Z., Yuan, Y., Sun, Y., Zhu, Y., Fu, Y., Jiang, Y., et al. (2024).
    Quantum classical hybrid convolutional neural networks for breast cancer diagnosis.
    \textit{Scientific Reports}, 14(1), 24699.
    \url{https://doi.org/10.1038/s41598-024-74778-7}

  % --- Referencia 16: Quantum Transfer Learning for Breast Cancer Detection ---

\bibitem{Zhang2013NEQR}
    Zhang, Y., Lu, K., Gao, Y., \& Wang, M. (2013).
    NEQR: A novel enhanced quantum representation of digital images.
    \textit{Quantum Information Processing}, 12(8), 2833--2860.
    \url{https://doi.org/10.1007/s11128-013-0567-z}

\bibitem{huang2021power}
    Huang, H.-Y., Broughton, M., Mohseni, M., Babbush, R., Boixo, S., Neven, H., \& McClean, J. R. (2021). Power of data in quantum machine learning.
    \textit{Nature Communications}, 12, 2631.
    \url{https://doi.org/10.1038/s41467-021-22539-9}

  % --- Referencia 44: Realization of Contrast Limited Adaptive Histogram Equalization ---

\bibitem{vanGriethuysen2017}
    van Griethuysen, J. J. M., et al. (2017). Computational radiomics system to decode the radiographic phenotype.
    \textit{Cancer Research}, 77(21), e104--e107.
    \url{https://doi.org/10.1158/0008-5472.CAN-17-0339}

  % --- Referencia 46: IBM Quantum Learning - Basics of Quantum Information ---
\end{thebibliography}
